% Chapter 1: scope, conventions and design principles

\chapter{Scope and conventions}

\label{Chapter1}

%----------------------------------------------------------------------------------------

\section{What this part covers}

This guide grows with the code. Part~I covers the two modules of \code{fluxlb.core} that
are complete and tested, the device and precision utilities in \file{gpu.py} and the
velocity sets in \file{lattices.py}, together with the tests that verify them. It stops
where the solver does: streaming, collision, boundary conditions and the time loop are not
yet implemented and are not described here.

\section{Authorship}

The repository's \file{CLAUDE.md} sets a three-tier ownership framework, summarised in
\cref{tab:tiers}. Scientific reasoning (Tier~1) is written by the maintainer; implementation
notes, tables and tooling (Tier~3) may be drafted from the code. In this document, orange
boxes mark Tier~1 passages still to be written, and green boxes record facts that were
checked numerically against the code as it stands.

\begin{table}[h]
\centering\small
\begin{tabular}{@{}lp{0.62\linewidth}@{}}
\toprule
Tier & Scope \\
\midrule
1 & Scientific logic and science prose: loss and residual formulations, boundary-condition
    choices, collision-operator and architecture rationale, physical reasoning, theory text. \\
2 & Drafting from direction: correspondence, changelogs, release notes. \\
3 & Execution to specification: boilerplate, plumbing, refactors, tooling, tests, CI. \\
\bottomrule
\end{tabular}
\caption{Intellectual-ownership tiers from \file{CLAUDE.md}.}
\label{tab:tiers}
\end{table}

\section{Notation}

A lattice has $Q$ discrete velocities $\vc_i$, $i = 0,\dots,Q-1$, in $D$ dimensions, with
weights $w_i$ and squared lattice sound speed $\cs^2$. The distribution tensor is $f$ with
shape $[Q, *\text{spatial}]$, so that index~0 runs over velocities. The opposite direction
of $i$ is $\bar{i}$, with $\vc_{\bar{i}} = -\vc_i$. Greek indices label spatial components
and $\delta_{\alpha\beta}$ is the Kronecker delta. The weighted lattice sum is
$\langle g \rangle = \sum_i w_i\, g(\vc_i)$.

%----------------------------------------------------------------------------------------

\section{Design principles carried by these modules}

Two conventions from \file{CLAUDE.md} shape everything in this part.

\begin{description}[style=nextline]
\item[One device and dtype context.]
Lattice constants are registered as PyTorch buffers so that a single
\code{module.to(device, dtype)} moves the whole solver. Nothing in a hot loop may query the
device or copy data to the host.

\item[fp32 compute, fp64 accumulation.]
Populations and collisions are computed in fp32. Conserved moments (density and momentum)
are accumulated in fp64 where the hardware allows it. The primary node has weak fp64
throughput and Apple's MPS backend has no fp64 at all, so the accumulation dtype is a
per-device decision made in one place (\cref{sec:precision}).
\end{description}

\begin{maintainer}[why these two conventions follow from the physics and the hardware]
Explain why conservation of mass and momentum, rather than the whole population update, is
the quantity that needs extended precision, and what error the fp32 fallback on MPS
introduces into the conservation checks.
\end{maintainer}
